\subsection{Definição dos objetivos (SMART)}

Para viabilizar a estratégia desenhada na visão de futuro da organização, estabelecem-se os seguintes objetivos de curto, médio e longo prazo, estruturados sob a metodologia SMART:

\subsubsection{Objetivos de Curto e Médio Prazo (Anos 1 a 3)}

\begin{itemize}
    \item \textbf{Adoção e Penetração de Mercado:} Alcançar a marca de 10.000 downloads ativos das ferramentas de software e projetos de hardware IoT até o final do segundo ano de operação, validando a usabilidade da plataforma no mercado de segurança cibernética e tomada de decisão.
    \item \textbf{Criação de Comunidade:} Consolidar uma comunidade ativa de desenvolvedores e criadores de conteúdo com, pelo menos, 500 membros engajados em fóruns oficiais e produção de documentação técnica até o final do terceiro ano.
    \item \textbf{Homologação do Produto:} Desenvolver, testar e colocar em ambiente de produção estável 3 variações de dispositivos de hardware IoT focados em acionamento resiliente de alertas até o décimo oitavo mês de atividade da empresa.
\end{itemize}

\subsubsection{Objetivos de Longo Prazo (Anos 4 e 5)}

\begin{itemize}
    \item \textbf{Monetização de Grandes Clientes:} Converter a ampla adoção do ecossistema em receita, conquistando uma carteira de, no mínimo, 50 grandes corporações integradas ao modelo de licenciamento pago até o final do quarto ano.
    \item \textbf{Liderança e Referência de Mercado:} Atingir o faturamento bruto anual de R\$ 5 milhões e estabelecer a marca como a maior referência nacional em sistemas de comunicação estratégica para mitigação de riscos cibernéticos até o encerramento do quinto ano de operação.
\end{itemize}
